\section{Open Questions (5)}

\subsection*{Q1. Passage-number effect on the dose-rate effect (DRE)}
\textbf{Question.} Does the dose-rate effect on G$_2$/M arrest in hADSCs persist across passages,
or does DRE magnitude decay with senescence-driven passage number?

\textbf{Basis.} The paper uses ATCC PCS-500-011 primary hADSCs but never reports the passage number,
and it is well-established that ADSC checkpoint plasticity and DNA-damage response kinetics change
with in-vitro age. With $n{=}3$ per condition and no passage variable, the 2021 study cannot separate
passage-conditioned response from dose-rate response.

\textbf{Next steps.} Repeat the identical LDR (1.40 Gy/min) vs HDR (7.31 Gy/min) at 2 Gy $^{137}$Cs
on the same ATCC lot at P3, P6, and P9 with $n{\ge}5$ per arm; score G$_2$/M at Day 1--3 and compute
per-passage effect size. Test whether the LDR-vs-HDR G$_2$/M gap regresses monotonically with passage.

\subsection*{Q2. ADSC-specific vs generic adherent-MSC response (ADSC vs BM-MSC)}
\textbf{Question.} Is the observed dose-rate sensitivity ADSC-specific, or a generic adherent-MSC /
adherent-culture response? Would matched bone-marrow-MSCs (BM-MSCs) show the same LDR/HDR G$_2$/M split?

\textbf{Basis.} The paper frames its findings as ADSC-specific but does not run a paired BM-MSC arm
on the same rig at the same rates. hADSC and hBM-MSC differ in glycolytic vs oxidative metabolism,
in p53 baseline, and in checkpoint stringency; the framing is under-supported without a lineage
comparator.

\textbf{Next steps.} Acquire ATCC PCS-500-012 hBM-MSCs (or equivalent), run the identical LDR/HDR
protocol side-by-side with hADSCs on the same $^{137}$Cs rig, and compare G$_2$/M drift magnitude
and CD44/TP53 fold changes head-to-head. Use a mixed-effects model to partition lineage vs dose-rate
variance.

\subsection*{Q3. Translation to iPSC-derived MSCs}
\textbf{Question.} Do iPSC-derived MSCs (which lack in-vivo aging and lack donor variance) recapitulate
the LDR-vs-HDR G$_2$/M pattern, and would that let us distinguish intrinsic checkpoint physics from
donor/passage confounds?

\textbf{Basis.} iPSC-derived MSCs are genetically defined, low-variance, and reproducible across labs.
If the LDR/HDR G$_2$/M split reproduces in iPSC-MSCs, the effect is mechanism-level; if it does not,
the 2021 paper's effect may be a primary-cell / donor artifact.

\textbf{Next steps.} Differentiate two iPSC lines (e.g.\ WTC-11, PGP1) to MSCs via published
trilineage protocols; irradiate at the same 2 Gy LDR/HDR; score G$_2$/M at Day 1--3 with the same
PI/Guava assay. Any effect-size ratio $<0.5\times$ the primary-hADSC value would flag a
donor/passage confound in the original paper.

\subsection*{Q4. Coupling to differentiation state}
\textbf{Question.} Is the G$_2$/M arrest coupled to differentiation state (undifferentiated vs
adipogenic-primed vs osteogenic-primed hADSCs)?

\textbf{Basis.} The paper studies undifferentiated hADSCs only. Adipogenic and osteogenic induction
reprogram cell-cycle regulation (CDKN1A, CCND1 dynamics) and could either amplify or abolish the
LDR/HDR G$_2$/M gap. The clinical relevance of ADSC dose-rate response depends critically on which
differentiation state is being irradiated in situ.

\textbf{Next steps.} Prime hADSCs for 48 h in adipogenic-induction medium and separately in
osteogenic-induction medium, then irradiate LDR vs HDR at 2 Gy. Score G$_2$/M by PI/Guava at Day 1--3,
and score lineage marker retention (PPARG for adipo, RUNX2 for osteo) by qPCR at Day 5 to test whether
dose-rate modulates lineage commitment.

\subsection*{Q5. In-vivo relevance for fat-graft radiotherapy}
\textbf{Question.} What is the in-vivo relevance for fat-graft radiotherapy? Do irradiated hADSCs
retain trophic and engraftment function after LDR vs HDR at clinically relevant fractional doses?

\textbf{Basis.} Fat-grafting post-lumpectomy or post-mastectomy delivers a bolus of ADSCs into a
tissue field that has often received fractionated radiotherapy (2 Gy/fx $\times$ 25--30). The 2021
paper's 2 Gy single-dose in-vitro model is a first step but does not test fractionation or the
trophic-factor secretome (VEGF, HGF, IL-6) that drives fat-graft take.

\textbf{Next steps.} Extend the 2 Gy single-dose to a 2 Gy $\times$ 5 fractionated regimen at LDR
and HDR; measure secretome (VEGF, HGF, IL-6, SDF-1) at 24 h post-final-fraction by multiplex ELISA;
and test engraftment in a NSG-mouse subcutaneous fat-graft model with irradiated vs sham hADSCs.
Score graft volume retention at 4 and 12 weeks.
